% !TEX program = xelatex
% 目标导向多模态情感分类的实证研究——基于 RethinkingTMSC 的中文转述
\documentclass[12pt]{ctexart}
\usepackage[a4paper,top=3.7cm,bottom=3.5cm,left=2.5cm,right=2.6cm]{geometry}
\usepackage{booktabs}
\usepackage{amsmath,amssymb}
\usepackage{graphicx}
\usepackage{caption}
\usepackage[colorlinks=true,linkcolor=blue,citecolor=blue,urlcolor=blue]{hyperref}

\linespread{1.5}
\emergencystretch=2em
\captionsetup{labelsep=quad, font={small}}

\ctexset{
  section={format={\heiti\zihao{3}\centering}, beforeskip=24pt, afterskip=12pt},
  subsection={format={\heiti\zihao{-4}}, beforeskip=12pt, afterskip=6pt},
}

\begin{document}
\thispagestyle{plain}

\begin{center}
{\heiti\zihao{-2} 目标导向多模态情感分类的实证研究}\par
\vspace{4pt}
{\zihao{-4} ——基于 RethinkingTMSC（EMNLP 2023 Findings）的中文转述}\par
\vspace{8pt}
{\zihao{5} 简相彪\quad 2026 年 8 月}
\end{center}

\vspace{10pt}
\noindent\textbf{摘要：}近年来，目标导向多模态情感分类（Target-oriented Multimodal
Sentiment Classification，TMSC）受到学者们的广泛关注，但现有多模态模型的性能已逐渐进入
瓶颈。为探究其成因，本文对主流模型结构与两个常用数据集开展系统的实证研究，围绕三个
问题展开：（1）各模态对 TMSC 是否同等重要？（2）哪类多模态融合模块更有效？（3）现有
数据集能否充分支撑该任务的研究？实验与数据分析表明：现有 TMSC 系统主要依赖文本模态，
多数目标的情感仅凭文本即可判定；各类融合模块相对纯文本模型没有显著增益，部分融合方式
甚至导致性能下降；数据集在规模、标注质量以及“图文共同决定情感”的样本占比等方面均存在
明显局限。基于上述发现，本文从模型设计与数据集构建两个角度给出若干建议，以期为 TMSC
的未来研究提供参考。

\noindent\textbf{关键词：}目标导向多模态情感分类；多模态融合；BERT；视觉编码器；实证研究

\clearpage
\section{引言}
目标导向情感分类又称方面级情感分类，是情感分析中的基础任务，其目标是判断文本中指定
目标实体所表达的情感极性（正面、中性或负面）\cite{pontiki2014}。为了进一步利用推文中
的图像信息，Yu 和 Jiang\cite{yu2019} 提出目标导向多模态情感分类（TMSC），将文本与图像
两种模态的信息相结合以判断目标的情感。

近年来，TMSC 系统的性能逐渐趋于平台期：在广泛使用的 Twitter15 与 Twitter17 两个数据集
上，现有最优基线模型的 Macro-F1 仅约 70\cite{ye2023}。为弄清性能停滞的原因，有必要从
模型层面与模态层面分别进行分析。粗略来看，现有模型结构可归结为两类模块：一是对各模态
进行表示学习的单模态编码器，二是建模模态间交互的多模态融合模块\cite{ye2023}。

本文在统一实验设置下系统比较不同编码器与融合模块，并对两个数据集的特性进行深入分析，
旨在回答以下三个问题：

\textbf{Q1}：各模态对 TMSC 是否同等重要？为回答该问题，本文分别评测文本与图像单模态
模型：文本侧采用预训练语言模型 BERT\cite{devlin2019}，图像侧采用 ResNet-152\cite{he2016}、
ViT\cite{dosovitskiy2021} 与 Faster R-CNN\cite{ren2015} 三种主流视觉编码器。

\textbf{Q2}：哪类多模态融合模块更有效？本文将现有融合策略归纳为拼接（Concatenation）、
张量融合（Tensor Fusion）\cite{zadeh2017}、自注意力（Self Attention）、图像到文本
（Image2Text）、文本到图像（Text2Image）与双向（Bi-direction）六类，并在统一设置下进行
对比。

\textbf{Q3}：现有数据集能否充分支撑该任务的研究？本文对 Twitter15 与 Twitter17 进行深入
分析，得到四点发现：（1）数据集规模有限；（2）多模态情感与文本情感的一致性远高于其与
视觉情感的一致性；（3）大量目标并未出现在图像中；（4）仅少数样本的情感需要文本与图像
共同决定\cite{ye2023}。

本文的主要贡献包括：（1）系统考察了不同单模态编码器与多模态融合模块对 TMSC 性能的
影响；（2）深入分析了现有数据集的不足；（3）基于实验观察给出若干研究建议，为 TMSC 的
模型设计与数据集构建提供参考\cite{ye2023}。

\clearpage
\section{相关工作}
作为情感分析的子任务，TMSC 近年来受到广泛关注\cite{yu2019}。Xu 等\cite{xu2019} 构建了
中文数据集 Multi-ZOL，并提出多跳记忆网络以建模模态间交互；随后 Yu 和 Jiang\cite{yu2019}
构建了 Twitter15 与 Twitter17 两个英文数据集，并将 BERT 引入该任务。后续研究大致沿两个
方向展开：一是不断探索增强模态间交互的方法\cite{khan2021}，二是将预训练模型应用于该
任务\cite{smp2022,vlp2022}。尽管已有大量工作，多模态模型相对纯文本模型仍未取得显著
增益，这正是本文开展实证研究的出发点。在文本建模方面，预训练语言模型 BERT\cite{devlin2019}
已取代早期基于 LSTM\cite{hochreiter1997} 与记忆网络\cite{weston2015} 的方法，成为该任务
的主流文本编码器。

\clearpage
\section{任务定义与模型体系}
\subsection{任务定义}
给定一条推文文本 $X$、文本中指定的目标实体 $t$ 以及推文配图 $I$，TMSC 的目标是判断目标
$t$ 被表达的情感极性 $y\in\{\text{正面，中性，负面}\}$。评测指标为准确率（ACC）与宏平均
F1（Macro-F1）\cite{ye2023}。

\subsection{单模态编码器}
文本编码器采用 BERT\cite{devlin2019}。对于图像模态，本文考察三种主流编码器：

（1）ResNet-152\cite{he2016}：将图像缩放为 $224\times224$ 后输入网络，取最后一层卷积输出
的 49 个区域作为图像表示，如公式~\eqref{eq:resnet} 所示：
\begin{equation}
  I=[v_1,v_2,\ldots,v_{49}],\quad v_i\in\mathbb{R}^{2048}
  \label{eq:resnet}
\end{equation}
其中 $v_i$ 表示第 $i$ 个区域的 2048 维特征向量。

（2）ViT\cite{dosovitskiy2021}：将图像划分为 $16\times16$ 的 patch，并在序列前端加入分类
标记（CLS token），经 Transformer\cite{vaswani2017} 编码后得到图像表示，如公式~\eqref{eq:vit}
所示：
\begin{equation}
  I=[v_{\mathrm{cls}},v_1,v_2,\ldots,v_{196}],\quad v_i\in\mathbb{R}^{768}
  \label{eq:vit}
\end{equation}

（3）Faster R-CNN\cite{ren2015}：采用在 Visual Genome\cite{krishna2017} 上重新训练的目标
检测模型，选取置信度最高的 36 个目标提议（object proposals）作为图像表示，如公式
~\eqref{eq:faster} 所示：
\begin{equation}
  I=[v_1,v_2,\ldots,v_{36}],\quad v_i\in\mathbb{R}^{2048}
  \label{eq:faster}
\end{equation}
其中 $v_i$ 来自区域提议网络（RPN）的 ROI 池化层。为统一不同图像编码器的输出维度，融合前
通过一个线性映射层将图像表示映射到 768 维\cite{ye2023}。

\subsection{多模态融合模块}
设池化后的文本表示为 $H_p^T\in\mathbb{R}^{768}$，池化后的图像表示为 $H_p^I\in\mathbb{R}^{768}$，
本文按以下六类融合策略进行对比\cite{ye2023}：

（1）拼接（Concatenate）：将文本与图像池化表示直接拼接，得到多模态表示（如公式
~\eqref{eq:concat} 所示）：
\begin{equation}
  H=H_p^I\oplus H_p^T,\quad H\in\mathbb{R}^{768+768}
  \label{eq:concat}
\end{equation}

（2）张量融合（Tensor Fusion）\cite{zadeh2017}：通过外积建模模态间的交互，同时保留各模态
自身的特性（如公式~\eqref{eq:tensor} 所示）：
\begin{equation}
  H=H_p^I\otimes H_p^T,\quad H\in\mathbb{R}^{768\times768}
  \label{eq:tensor}
\end{equation}

（3）自注意力（Self Attention）：将图像表示与文本表示拼接后，经过三层自注意力层与池化层，
得到（如公式~\eqref{eq:selfattn} 所示）：
\begin{equation}
  H=\mathrm{SelfAttn}([H^T;H^I])\in\mathbb{R}^{768}
  \label{eq:selfattn}
\end{equation}

（4）图像到文本（Image2Text）与文本到图像（Text2Image）：前者以图像表示作为查询
（Query）、文本表示作为键值（Key/Value），经三层交叉注意力得到融合表示；后者将两者角色
互换。将二者结果拼接，即得到双向（Bi-direction）表示（如公式~\eqref{eq:bi} 所示）
\cite{ye2023}：
\begin{equation}
  H=[H^{I\rightarrow T};H^{T\rightarrow I}]\in\mathbb{R}^{768+768}
  \label{eq:bi}
\end{equation}
其中 $H^{I\rightarrow T}$ 与 $H^{T\rightarrow I}$ 分别表示 Image2Text 与 Text2Image 的输出。
表~\ref{tab:models} 汇总了本文涉及的各类基线模型的结构差异。

\begin{table}[htbp]
\centering
\caption{各类基准模型的结构总览\cite{ye2023}}
\label{tab:models}
\small
\begin{tabular}{lccc}
\toprule
模型 & 文本编码器 & 图像编码器 & 融合策略 \\
\midrule
mBERT\cite{yu2019} & BERT & ResNet-152 & 拼接 \\
TomBERT\cite{yu2019} & BERT & ResNet-152 & 拼接 + 图像到文本 \\
EF-CapTrBERT\cite{khan2021} & BERT & ResNet-152 & 跨模态交互 \\
Res-BERT+BL\cite{ye2023} & BERT & ResNet-152 & 自注意力 \\
Res-BERT+BL-TFN\cite{ye2023} & BERT & ResNet-152 & 张量融合 \\
SMP\cite{smp2022} & BERT & ViT & 双向注意力 \\
VLP\cite{vlp2022} & 预训练视觉语言模型 & Faster R-CNN & 预训练 \\
\bottomrule
\end{tabular}
\end{table}

\clearpage
\section{实验设置与结果分析}
\subsection{实验设置}
实验在 Twitter15 与 Twitter17 两个英文数据集上进行\cite{yu2019}。统一超参数设置如下：
训练 8 个 epoch，批大小 32，学习率 $2\times10^{-5}$，优化器采用 Adam\cite{kingma2015}，
并在 5 个随机种子（0/42/199/2022/11122）上重复实验，报告均值与标准差\cite{ye2023}。

\subsection{总体结果}
表~\ref{tab:results} 给出各单模态模型与多模态模型在两个数据集上的实验结果。整体来看，
可以得出以下四点观察\cite{ye2023}：

第一，纯文本模型（BERT）表现良好，而纯视觉模型（ResNet、ViT、Faster R-CNN）表现较差，
说明在该任务上模型对文本的依赖远大于图像，且这一现象在 Twitter15 上更为明显。

第二，融合方式对性能有显著影响：以获取文本信息为主的融合（如 Image2Text）优于以获取
图像信息为主的融合（如 Text2Image），再次印证文本与图像在任务中的重要性并不对等。

第三，与纯文本模型相比，各类融合模块并未带来显著增益，部分融合甚至更差。其原因在于
部分图像并未提供相关信息，反而引入了干扰信息。

第四，图像编码器的影响并不明确：单模态图像模型性能低且方差大，而在多模态融合设置下，
不同图像编码器之间的性能差异很小。这与现有数据集中视觉数据的特性有关，详见第 5 节的
数据分析。

基于上述分析，未来 TMSC 模型设计应重点关注：（1）充分利用文本信息的优势；（2）设计
更有效的图像编码方法，以更好地提取图像语义信息；（3）增强融合模块的抗噪能力，使其能够
灵活选择并利用文本与图像中的有效特征\cite{ye2023}。

\begin{table}[htbp]
\centering
\caption{Twitter15 与 Twitter17 上的实验结果（5 种子均值±标准差）\cite{ye2023}}
\label{tab:results}
\scriptsize
\resizebox{\textwidth}{!}{%
\begin{tabular}{lcccccc}
\toprule
模态 & 模型 & 融合 & T15 ACC & T15 F1 & T17 ACC & T17 F1 \\
\midrule
文本 & BERT & --- & 76.72±1.16 & 71.19±2.19 & 68.04±0.40 & 65.66±0.35 \\
图像 & ResNet & --- & 57.65±1.00 & 32.52±2.66 & 57.79±0.99 & 51.98±1.23 \\
图像 & ViT & --- & 59.65±1.13 & 31.25±2.71 & 59.53±0.95 & 54.08±0.78 \\
图像 & Faster R-CNN & --- & 55.97±1.10 & 35.72±5.43 & 56.18±0.85 & 49.88±1.70 \\
多模态 & ResNet & 拼接 & 75.29±0.45 & 68.71±1.34 & 67.92±0.56 & 65.32±0.53 \\
多模态 & ResNet & 张量融合 & 74.19±0.94 & 68.93±0.57 & 66.66±1.21 & 63.99±1.61 \\
多模态 & ResNet & 自注意力 & 76.03±0.96 & 70.57±2.39 & 68.01±0.96 & 65.41±1.60 \\
多模态 & ResNet & Image2Text & 77.13±1.33 & 71.48±1.90 & 69.37±0.36 & 66.85±0.79 \\
多模态 & ResNet & Text2Image & 75.18±1.66 & 67.77±4.81 & 68.07±0.58 & 65.18±1.48 \\
多模态 & ResNet & 双向 & 77.32±0.63 & 72.06±0.81 & 68.41±1.01 & 66.39±1.39 \\
多模态 & ViT & 拼接 & 76.22±0.90 & 70.37±1.45 & 67.94±0.70 & 66.17±0.78 \\
多模态 & ViT & 张量融合 & 73.44±0.78 & 67.46±1.45 & 65.46±1.67 & 62.02±1.40 \\
多模态 & ViT & 自注意力 & 75.08±0.41 & 68.94±0.83 & 67.52±0.58 & 65.56±0.35 \\
多模态 & ViT & Image2Text & 77.11±0.44 & 71.91±0.42 & 69.14±0.52 & 66.96±0.68 \\
多模态 & ViT & Text2Image & 75.12±1.01 & 69.40±1.38 & 67.52±1.06 & 64.49±1.46 \\
多模态 & ViT & 双向 & 76.70±0.75 & 71.67±1.45 & 69.16±0.17 & 67.25±0.56 \\
多模态 & Faster & 拼接 & 75.45±0.73 & 69.77±1.23 & 67.60±1.15 & 64.74±1.69 \\
多模态 & Faster & 张量融合 & 72.09±0.66 & 66.77±1.04 & 66.34±1.45 & 62.96±2.09 \\
多模态 & Faster & 自注意力 & 76.09±0.89 & 70.08±1.37 & 68.09±1.10 & 66.12±1.23 \\
多模态 & Faster & Image2Text & 77.36±0.37 & 71.69±0.37 & 68.43±0.65 & 66.44±1.10 \\
多模态 & Faster & Text2Image & 70.82±2.99 & 57.94±5.81 & 60.31±6.43 & 54.50±7.06 \\
多模态 & Faster & 双向 & 76.57±0.46 & 70.88±0.89 & 69.51±0.62 & 67.50±0.37 \\
\bottomrule
\end{tabular}}
\end{table}

\clearpage
\section{数据集分析}
\subsection{数据集统计}
表~\ref{tab:data} 给出两个数据集的统计信息。可以看出：（1）数据集规模较小，每个样本平均
包含的目标数不足 1.5 个；（2）情感标签分布不均衡，中性样本约占 50\%，负面样本不足 15\%。
造成上述现象的原因在于，Twitter15 与 Twitter17 最初分别是为命名实体识别任务构建的
\cite{zhang2018,lu2018}，而非专门为 TMSC 设计\cite{ye2023}。

\begin{table}[htbp]
\centering
\caption{两个数据集的统计信息（“均目标”指每个样本的平均目标数）\cite{ye2023}}
\label{tab:data}
\scriptsize
\resizebox{\textwidth}{!}{%
\begin{tabular}{lcccccccccc}
\toprule
划分 & T15 负面 & T15 中性 & T15 正面 & T15 总数 & T15 均目标 &
T17 负面 & T17 中性 & T17 正面 & T17 总数 & T17 均目标 \\
\midrule
训练 & 368 & 1883 & 928 & 3179 & 1.348 & 416 & 1638 & 1508 & 3562 & 1.410 \\
开发 & 149 & 670 & 303 & 1122 & 1.336 & 144 & 517 & 515 & 1176 & 1.439 \\
测试 & 113 & 607 & 317 & 1037 & 1.354 & 168 & 573 & 493 & 1234 & 1.450 \\
\bottomrule
\end{tabular}}
\end{table}

\subsection{标注层面分析}
为进一步理解性能瓶颈，本文参照 Yu 和 Jiang\cite{yu2019} 的标注流程，邀请三位领域专家
从 Twitter15 与 Twitter17 测试集中各随机抽取 200 条样本，从文本情感一致性、视觉情感
一致性、目标是否存在以及情感是否由图文共同决定四个方面进行标注，并采用多数投票作为
最终标注结果（见表~\ref{tab:anno}）\cite{ye2023}。分析得到以下发现：

第一，多模态情感与文本情感具有高度一致性，但与视觉情感的一致性较低。在 Twitter15 上，
93\% 的目标其文本情感与多模态情感一致，而视觉情感一致的比例仅为 47.5\%。这说明数据集
存在明显的分布偏差，即文本信息对判定多模态情感更为有效。该现象在 Twitter17 上有所缓解，
但文本信息的一致性仍高于视觉信息。

第二，大量目标并未出现在图像中，这不符合目标导向多模态情感分类任务的设定。其原因可能
在于数据集构建时目标直接从文本中选取，并未考虑对应图像的内容\cite{yu2019}。

第三，由于图像与目标不相关以及目标在图像中缺失，仅有少量样本的情感需要文本与图像共同
决定：Twitter15 上约为 22\%，Twitter17 上约为 55\%。就多模态任务而言，这两个数据集可能
并非最合适的评测基准\cite{ye2023}。

\begin{table}[htbp]
\centering
\caption{400 条随机样本的人工标注分析（三位专家多数投票）\cite{ye2023}}
\label{tab:anno}
\small
\begin{tabular}{lcc}
\toprule
标注指标 & Twitter15 & Twitter17 \\
\midrule
文本情感与多模态情感一致 & 93\% & 较高（仍高于视觉一致性） \\
视觉情感与多模态情感一致 & 47.5\% & 较低（较 Twitter15 有所缓解） \\
情感需图文共同决定 & 22\% & 55\% \\
\bottomrule
\end{tabular}
\end{table}

基于上述分析，本文认为高质量 TMSC 数据集应具备以下特征：（1）真实反映现实世界的数据
分布，包括标签不均衡等情况，并为不同情形提供充足的样本；（2）具有较大的数据多样性，
涵盖多种数据类型与领域，以有效检验模型的泛化能力与鲁棒性；（3）提供多维度标注信息，
包括多模态情感与单模态情感，以支持对模型处理不同数据来源能力的深入分析\cite{ye2023}。

\clearpage
\section{结论与展望}
本文对 TMSC 任务开展了系统的实证研究与数据分析。结果表明：现有多模态模型相对纯文本
模型并未取得显著性能提升，其主要原因在于现有数据集对文本模态的过度依赖，视觉模态发挥
的作用相对有限。基于实验分析，本文从模型设计与数据集构建两个方面提出了未来研究方向，
以更好地体现社交媒体情感的多模态特性。

本文工作仍存在一定局限：（1）数据分析主要针对当前公开的英文数据集 Twitter15 与
Twitter17，未涉及研究较少的中文数据集 Multi-ZOL\cite{xu2019}；（2）虽然分析指出了现有
数据集的问题，但并未构建新的、更合适的数据集；（3）实验未专门比较不同文本编码方法的
影响，而将研究重点放在图像编码方法与融合模块上\cite{ye2023}。上述问题将作为后续工作
继续探索。

\clearpage
\begin{thebibliography}{19}
\small
\bibitem{ye2023} Ye J, Zhou J, Tian J, et al. RethinkingTMSC: An Empirical Study for
Target-Oriented Multimodal Sentiment Classification[C]. Findings of ACL: EMNLP 2023: 270--277.
\bibitem{yu2019} Yu J, Jiang J. Adapting BERT for Target-Oriented Multimodal Sentiment
Classification[C]. IJCAI 2019: 5408--5414.
\bibitem{devlin2019} Devlin J, Chang M W, Lee K, et al. BERT: Pre-training of Deep Bidirectional
Transformers for Language Understanding[C]. NAACL-HLT 2019: 4171--4186.
\bibitem{he2016} He K, Zhang X, Ren S, et al. Deep Residual Learning for Image Recognition[C].
CVPR 2016: 770--778.
\bibitem{dosovitskiy2021} Dosovitskiy A, Beyer L, Kolesnikov A, et al. An Image is Worth
$16\times16$ Words: Transformers for Image Recognition at Scale[C]. ICLR 2021.
\bibitem{ren2015} Ren S, He K, Girshick R, et al. Faster R-CNN: Towards Real-Time Object
Detection with Region Proposal Networks[C]. NeurIPS 2015: 91--99.
\bibitem{zadeh2017} Zadeh A, Chen M, Poria S, et al. Tensor Fusion Network for Multimodal
Sentiment Analysis[C]. EMNLP 2017: 1103--1114.
\bibitem{vaswani2017} Vaswani A, Shazeer N, Parmar N, et al. Attention is All You Need[C].
NeurIPS 2017: 5998--6008.
\bibitem{zhang2018} Zhang Q, Fu J, Liu X, et al. Adaptive Co-attention Network for Named
Entity Recognition in Tweets[C]. AAAI 2018: 5674--5681.
\bibitem{lu2018} Lu D, Neves L, Carvalho V, et al. Visual Attention Model for Name Tagging in
Multimodal Social Media[C]. ACL 2018: 1990--1999.
\bibitem{pontiki2014} Pontiki M, Galanis D, Pavlopoulos J, et al. SemEval-2014 Task 4: Aspect
Based Sentiment Analysis[C]. SemEval@COLING 2014: 27--35.
\bibitem{hochreiter1997} Hochreiter S, Schmidhuber J. Long Short-Term Memory[J]. Neural
Computation, 1997, 9(8): 1735--1780.
\bibitem{weston2015} Weston J, Chopra S, Bordes A. Memory Networks[C]. ICLR 2015.
\bibitem{khan2021} Khan Z, Fu Y. Exploiting BERT for Multimodal Target Sentiment Classification
through Input Space Translation[C]. ACM MM 2021: 3034--3042.
\bibitem{smp2022} Ye J, Zhou J, Tian J, et al. Sentiment-Aware Multimodal Pre-training for
Multimodal Sentiment Analysis[J]. Knowledge-Based Systems, 2022, 258: 110021.
\bibitem{vlp2022} Ling Y, Yu J, Xia R. Vision-Language Pre-training for Multimodal Aspect-Based
Sentiment Analysis[C]. ACL 2022: 2149--2159.
\bibitem{krishna2017} Krishna R, Zhu Y, Groth O, et al. Visual Genome: Connecting Language and
Vision Using Crowdsourced Dense Image Annotations[J]. IJCV, 2017, 123(1): 32--73.
\bibitem{kingma2015} Kingma D P, Ba J. Adam: A Method for Stochastic Optimization[C]. ICLR 2015.
\bibitem{xu2019} Xu N, Mao W, Chen G. Multi-interactive Memory Network for Aspect Based
Multimodal Sentiment Analysis[C]. AAAI 2019: 371--378.
\end{thebibliography}

\end{document}
